\documentclass[10pt,a4paper]{amsart}
\usepackage{amsmath,amssymb,amsthm,graphicx}
% Retained preamble for the Stacks Project chapter snapshot in this directory.
% The chapter's own preamble.tex is a copy of the shared Stacks Project preamble; the block
% below is sliced MECHANICALLY, line for line, from it. Nothing was added, removed or
% redefined; the wrapper class, the xr-hyper/hyperref wiring and the externaldocument list
% are omitted because the assembler supplies the document class.
%
% lines taken from preamble.tex: 3, 5, 9, 11, 12, 156, 157, 158, ... (51 lines total)
%
% package coverage restored from the chapter's own preamble (2026-10-13): xy, xy-2cell

\usepackage[all]{xy}
\xyoption{2cell}
\UseAllTwocells

\newenvironment{reference}{\comment}{\endcomment}
\newenvironment{slogan}{\comment}{\endcomment}
\newenvironment{history}{\comment}{\endcomment}
\theoremstyle{plain}
\newtheorem{theorem}[subsection]{Theorem}
\newtheorem{proposition}[subsection]{Proposition}
\newtheorem{lemma}[subsection]{Lemma}
\theoremstyle{definition}
\newtheorem{definition}[subsection]{Definition}
\newtheorem{example}[subsection]{Example}
\newtheorem{exercise}[subsection]{Exercise}
\newtheorem{situation}[subsection]{Situation}
\theoremstyle{remark}
\newtheorem{remark}[subsection]{Remark}
\newtheorem{remarks}[subsection]{Remarks}
\numberwithin{equation}{subsection}
\def\lim{\mathop{\mathrm{lim}}\nolimits}
\def\colim{\mathop{\mathrm{colim}}\nolimits}
\def\Spec{\mathop{\mathrm{Spec}}}
\def\Hom{\mathop{\mathrm{Hom}}\nolimits}
\def\Ext{\mathop{\mathrm{Ext}}\nolimits}
\def\SheafHom{\mathop{\mathcal{H}\!\mathit{om}}\nolimits}
\def\SheafExt{\mathop{\mathcal{E}\!\mathit{xt}}\nolimits}
\def\Sch{\mathit{Sch}}
\def\Mor{\mathop{\mathrm{Mor}}\nolimits}
\def\Ob{\mathop{\mathrm{Ob}}\nolimits}
\def\Sh{\mathop{\mathit{Sh}}\nolimits}
\def\NL{\mathop{N\!L}\nolimits}
\def\CH{\mathop{\mathrm{CH}}\nolimits}
\def\proetale{{pro\text{-}\acute{e}tale}}
\def\etale{{\acute{e}tale}}
\def\QCoh{\mathit{QCoh}}
\def\Ker{\mathop{\mathrm{Ker}}}
\def\Im{\mathop{\mathrm{Im}}}
\def\Coker{\mathop{\mathrm{Coker}}}
\def\Coim{\mathop{\mathrm{Coim}}}
\DeclareMathSymbol{\boxtimes}{\mathbin}{AMSa}{"02}
\def\QCohstack{\mathcal{QC}\!\mathit{oh}}
\def\Cohstack{\mathcal{C}\!\mathit{oh}}
\def\Spacesstack{\mathcal{S}\!\mathit{paces}}
\def\Quotfunctor{\mathrm{Quot}}
\def\Hilbfunctor{\mathrm{Hilb}}
\def\Curvesstack{\mathcal{C}\!\mathit{urves}}
\def\Polarizedstack{\mathcal{P}\!\mathit{olarized}}
\def\Complexesstack{\mathcal{C}\!\mathit{omplexes}}
\def\Pic{\mathop{\mathrm{Pic}}\nolimits}
\def\Picardstack{\mathcal{P}\!\mathit{ic}}
\def\Picardfunctor{\mathrm{Pic}}
\def\Deformationcategory{\mathcal{D}\!\mathit{ef}}
% lane-s11: discard bibliographic \begin{reference} blocks entirely (incl. any \cite inside).
\newsavebox{\limrefbox}
\ifcsname reference\endcsname\expandafter\let\csname reference\endcsname\undefined\fi\ifcsname endreference\endcsname\expandafter\let\csname endreference\endcsname\undefined\fi
\def\citation#1{\relax}
\newenvironment{reference}{\setbox\limrefbox=\vbox\bgroup\let\cite\relax}{\egroup\ignorespaces}
\usepackage{cleveref}
\crefname{originallemma}{original lemma}{original lemmas}
\Crefname{originallemma}{Original lemma}{Original lemmas}
\makeatletter
\newlabel{lemma-descend-modules-finite-presentation}{{01ZR}{source}{}{sourcepointer-1}{}}
\newlabel{lemma-descend-modules-finite-presentation@cref}{{[originallemma][0][]01ZR}{[1][0][]source}{}{}{}}
\makeatother
\crefname{originalreference}{original reference}{original references}
\Crefname{originalreference}{Original reference}{Original references}
\makeatletter
\newlabel{algebra-lemma-flat-finite-presentation-limit-flat}{{02JO}{source}{}{sourcepointer-2}{}}
\newlabel{algebra-lemma-flat-finite-presentation-limit-flat@cref}{{[originalreference][0][]02JO}{[1][0][]source}{}{}{}}
\makeatother
\crefname{originalreference}{original reference}{original references}
\Crefname{originalreference}{Original reference}{Original references}
\makeatletter
\newlabel{algebra-lemma-finite-projective}{{00NX}{source}{}{sourcepointer-3}{}}
\newlabel{algebra-lemma-finite-projective@cref}{{[originalreference][0][]00NX}{[1][0][]source}{}{}{}}
\makeatother
\crefname{originalreference}{original reference}{original references}
\Crefname{originalreference}{Original reference}{Original references}
\makeatletter
\newlabel{modules-lemma-invertible}{{0B8K}{source}{}{sourcepointer-4}{}}
\newlabel{modules-lemma-invertible@cref}{{[originalreference][0][]0B8K}{[1][0][]source}{}{}{}}
\makeatother
\crefname{originalreference}{original reference}{original references}
\Crefname{originalreference}{Original reference}{Original references}
\makeatletter
\newlabel{algebra-lemma-module-map-property-in-colimit}{{05LI}{source}{}{sourcepointer-5}{}}
\newlabel{algebra-lemma-module-map-property-in-colimit@cref}{{[originalreference][0][]05LI}{[1][0][]source}{}{}{}}
\makeatother
\title{Descent of invertible modules along limits of schemes}
\author{The Stacks Project Authors}
\date{stacks-project:limits.tex:local-snapshot-20261012; GNU FDL 1.2 or later}
\begin{document}
\maketitle
\noindent\textit{Editorial note.} Author omission kept verbatim: 'After increasing \$i\$ we may assume this map is surjective (details omitted; hint: use Algebra, Lemma \textbackslash{}ref\{algebra-lemma-module-map-property-in-colimit\} on affine open cover)'. The omission is auxiliary (a lift existence step) and carries a hint plus a named proved lemma.\par\medskip
\noindent\textit{Allowed standard prerequisites (editorial).}\begin{itemize}
\item Module map property in colimit. Algebra, Lemma algebra-lemma-module-map-property-in-colimit
\end{itemize}\medskip
\ifdefined\hypertarget\hypertarget{sourcepointer-1}{}\fi
\noindent\textit{Source reference.} Original lemma 01ZR: sources/stacks-chow/tags.txt, source lines 2604--2604. Cross-chapter standard prerequisite cited by the author inside the counted span; it is used as a pointer and its own proof is not part of this item.\par\medskip
\ifdefined\hypertarget\hypertarget{sourcepointer-2}{}\fi
\noindent\textit{Source reference.} Original reference 02JO: sources/stacks-chow/tags.txt, source lines 3343--3343. Cross-chapter standard prerequisite cited by the author inside the counted span; it is used as a pointer and its own proof is not part of this item.\par\medskip
\ifdefined\hypertarget\hypertarget{sourcepointer-3}{}\fi
\noindent\textit{Source reference.} Original reference 00NX: sources/stacks-chow/tags.txt, source lines 875--875. Cross-chapter standard prerequisite cited by the author inside the counted span; it is used as a pointer and its own proof is not part of this item.\par\medskip
\ifdefined\hypertarget\hypertarget{sourcepointer-4}{}\fi
\noindent\textit{Source reference.} Original reference 0B8K: sources/stacks-chow/tags.txt, source lines 14143--14143. Cross-chapter standard prerequisite cited by the author inside the counted span; it is used as a pointer and its own proof is not part of this item.\par\medskip
\ifdefined\hypertarget\hypertarget{sourcepointer-5}{}\fi
\noindent\textit{Source reference.} Original reference 05LI: sources/stacks-chow/tags.txt, source lines 7237--7237. Cross-chapter standard prerequisite cited by the author inside the counted span; it is used as a pointer and its own proof is not part of this item.\par\medskip
\begin{lemma}
\label{lemma-descend-invertible-modules}
Let $S = \lim S_i$ be the limit of a directed system of quasi-compact and
quasi-separated schemes $S_i$ with affine transition morphisms. Then
\begin{enumerate}
\item any finite locally free $\mathcal{O}_S$-module is the pullback
of a finite locally free $\mathcal{O}_{S_i}$-module for some $i$,
\item any invertible $\mathcal{O}_S$-module is the pullback of an invertible
$\mathcal{O}_{S_i}$-module for some $i$, and
\item any finite type quasi-coherent ideal $\mathcal{I} \subset \mathcal{O}_S$
is of the form $\mathcal{I}_i \cdot \mathcal{O}_S$ for some $i$ and some
finite type quasi-coherent ideal $\mathcal{I}_i \subset \mathcal{O}_{S_i}$.
\end{enumerate}
\end{lemma}

\begin{proof}
Let $\mathcal{E}$ be a finite locally free $\mathcal{O}_S$-module. Since
finite locally free modules are of finite presentation we can find an $i$
and an $\mathcal{O}_{S_i}$-module $\mathcal{E}_i$ of finite presentation
such that $f_i^*\mathcal{E}_i \cong \mathcal{E}$, see
Lemma \ref{lemma-descend-modules-finite-presentation}.
After increasing $i$ we may assume $\mathcal{E}_i$ is a flat
$\mathcal{O}_{S_i}$-module, see
Algebra, Lemma \ref{algebra-lemma-flat-finite-presentation-limit-flat}.
(Using this lemma is not necessary, but it is convenient.)
Then $\mathcal{E}_i$ is finite locally free by
Algebra, Lemma \ref{algebra-lemma-finite-projective}.

\medskip\noindent
If $\mathcal{L}$ is an invertible $\mathcal{O}_S$-module,
then by the above we can find an $i$ and finite locally
free $\mathcal{O}_{S_i}$-modules $\mathcal{L}_i$
and $\mathcal{N}_i$ pulling back to $\mathcal{L}$ and
$\mathcal{L}^{\otimes -1}$. After possible increasing
$i$ we see that the map
$\mathcal{L} \otimes_{\mathcal{O}_X} \mathcal{L}^{\otimes -1}
\to \mathcal{O}_X$ descends to a map
$\mathcal{L}_i \otimes_{\mathcal{O}_{S_i}} \mathcal{N}_i \to
\mathcal{O}_{S_i}$. And after increasing $i$ further, we
may assume it is an isomorphism. It follows that
$\mathcal{L}_i$ is an invertible module
(Modules, Lemma \ref{modules-lemma-invertible}) and
the proof of (2) is complete.

\medskip\noindent
Given $\mathcal{I}$ as in (3) we see that
$\mathcal{O}_S \to \mathcal{O}_S/\mathcal{I}$
is a map of finitely presented $\mathcal{O}_S$-modules.
Hence by Lemma \ref{lemma-descend-modules-finite-presentation}
this is the pullback of some
map $\mathcal{O}_{S_i} \to \mathcal{F}_i$ of finitely
presented $\mathcal{O}_{S_i}$-modules. After increasing $i$
we may assume this map is surjective (details omitted; hint: use
Algebra, Lemma \ref{algebra-lemma-module-map-property-in-colimit}
on affine open cover). Then the kernel of $\mathcal{O}_{S_i} \to \mathcal{F}_i$
is a finite type quasi-coherent ideal in $\mathcal{O}_{S_i}$
whose pullback gives $\mathcal{I}$.
\end{proof}

\end{document}
